\chapter{Implementation in pystencils}\label{r02:ch:pystencils}

\begin{r02fork}
  All changes in this chapter are in commit \texttt{9a989e16} (``rvv support opus5.5
  VLS only VLEN256'') on the \texttt{rvv-support} branch of \file{submodules/pystencils},
  on top of upstream \texttt{20d7dcf0}. Report~1, \cref{r01:ch:rvv-backend}, documents
  the backend from a \ps{} user's perspective with more generated examples.
\end{r02fork}

From waLBerla's point of view, \ps{} is a black box that turns a target and a
configuration into C++ source. This chapter opens that box just far enough to show what
was added for RVV and how it plugs into the existing structure. Every new piece has an
exact counterpart in the x86 backend.

\section{Overview}

\begin{center}\small
\begin{tabularx}{\linewidth}{@{}>{\raggedright\arraybackslash}p{6.4cm}r>{\raggedright\arraybackslash}X@{}}
  \toprule
  \textbf{File (under \file{src/pystencils/})} & \textbf{Lines} & \textbf{Role / x86 counterpart}\\
  \midrule
  \file{codegen/target.py} & +43 & \code{Target.RISCV_RVV}, default lanes, host detection\\
  \file{codegen/driver.py} & +4 & maps the target to the platform\\
  \file{backend/platforms/rvv.py} & +490 (new) & platform and intrinsic selector;
    \file{x86.py}\\
  \file{backend/platforms/__init__.py} & +2 & export \cls{RvvCpu}\\
  \file{include/pystencils_runtime/rvv.hpp} & +189 (new) & runtime helpers;
    \file{x86.hpp}\\
  \file{jit/cpu/compiler_info.py} & +27 & \code{-march} for \code{RISCV_RVV} and for
    \code{CurrentCPU} on riscv64\\
  \midrule
  \file{tests/nbackend/test_rvv_{target,codegen,execution}.py} & +479 (new) & RVV unit,
    codegen and hardware tests\\
  six shared test modules & +40 & RVV cases next to AVX and SVE\\
  \bottomrule
\end{tabularx}
\end{center}
No IR node, transformation or loop strategy was changed. That is the practical meaning
of ``AVX parity'': the vectorizer produces the same vector IR for RVV as for AVX, and
only the last step, intrinsic selection, differs.

\section{The target}\label{r02:sec:ps-target}

\ps{} describes targets as combinations of feature flags in an \code{enum.Flag}. The new
target adds two architecture bits:
\begin{pycode}
class Target(Flag):
    ...
    _RISCV = auto()
    _RVV = auto()
    ...
    RISCV_RVV = _CPU | _VECTOR | _RISCV | _RVV
    """RISC-V architecture with the ratified V vector extension (RVV 1.0). Treated like a
    fixed-width target (cf. `X86_AVX`): VLEN = 256 bits, vector-length-specific code."""
\end{pycode}
\code{default_vector_lanes} returns $256/\text{element bits}$ for \code{RISCV_RVV},
the same case as \code{X86_AVX}: 4 lanes for \code{double}, 8 for \code{float}.

Host detection is the change that makes waLBerla's \code{CurrentCPU} work.
\code{_available_vector_targets()} gained a branch for riscv64 that calls the new helper
\code{_riscv_has_rvv()} (\cref{r02:lst:detect}); \code{Target.auto_cpu()} then returns
\code{RISCV_RVV} on the K1.

\rneedspace{12\baselineskip}
\begin{pycode}[nolol=false,caption={[RVV host detection in pystencils]RVV host detection in \texttt{codegen/target.py} (abridged).},label=r02:lst:detect]
    elif machine_spec == "riscv64":
        if _riscv_has_rvv():
            targets.append(Target.RISCV_RVV)
    ...

def _riscv_has_rvv(cpuinfo: str | None = None) -> bool:
    """RVV 1.0 shows up as the single-letter extension `v` in the leading part of the
    ISA string before the first `_`. `_zve64d` or `_xtheadvector` (RVV 0.7) must not count."""
    if cpuinfo is None:
        try:
            with open("/proc/cpuinfo") as f:
                cpuinfo = f.read()
        except OSError:
            return False
    for line in cpuinfo.splitlines():
        key, _, value = line.partition(":")
        if key.strip() == "isa":
            base = value.strip().lower().split("_")[0]           # e.g. "rv64imafdcv"
            return base.startswith("rv") and "v" in base[4:]     # skip "rv64"/"rv32"
    return False
\end{pycode}
The \code{cpuinfo} argument is a testing hook: \file{test_rvv_target.py} feeds it ISA
strings of real and invented chips. In \code{DefaultKernelCreationDriver._get_platform},
the target maps to the platform class:
\begin{pycode}
            elif self._target == Target.RISCV_RVV:
                from ..backend.platforms.rvv import RvvCpu
                return RvvCpu(self._ctx)
\end{pycode}

\section{The platform and its intrinsic selector}\label{r02:sec:ps-platform}

\file{backend/platforms/rvv.py} holds two classes, modelled on \file{x86.py}:
\begin{description}[style=nextline]
  \item[\cls{RvvCpu}(\cls{GenericVectorCpu})] the platform. It adds
    \code{"pystencils_runtime/rvv.hpp"} to the required headers and returns the
    intrinsic selector. It refuses \code{use_builtin_convertvector}, because RVV types
    are sizeless and cannot be used with GCC/Clang vector extensions.
  \item[\cls{SelectIntrinsicsRvv}(\cls{SelectIntrinsics})] the selector. The base class
    walks the vector IR and calls one hook per construct; this class implements the
    hooks.
\end{description}
Everything is driven by three small naming helpers. The register-group multiplier is
derived from the vector width exactly the way AVX picks between \code{__m256} and
\code{__m128}:
\rneedspace{12\baselineskip}
\begin{pycode}[nolol=false,caption={[Type naming in the RVV platform]Type naming in \texttt{backend/platforms/rvv.py}.},label=r02:lst:lmul]
RVV_VLEN = 256   # the single knob: fixing VLEN turns RVV into a fixed-width SIMD ISA

def _lmul(vtype: PsVectorType) -> str:
    match vtype.width:
        case w if w == RVV_VLEN:
            return "m1"                                   # AVX: __m256*
        case w if w == RVV_VLEN // 2 and vtype.scalar_type.width <= 32:
            return "mf2"                                  # AVX: __m128*; no vfloat64mf2_t
        case _:
            raise MaterializationError(...)               # wider vectors: rejected, like AVX

def _sfx(vtype):      # "f64m1", "i32mf2", ...  -> intrinsic name suffix
def _rvv_type(vtype): # PsCustomType("vfloat64m1_t"), tuples "vfloat32m1x4_t"
\end{pycode}

The hooks then assemble intrinsic names from a stem and the suffix:
\begin{center}\small
\begin{tabularx}{\linewidth}{@{}>{\raggedright\arraybackslash}p{3.4cm}>{\raggedright\arraybackslash}X@{}}
  \toprule
  \textbf{Hook} & \textbf{Generated code}\\
  \midrule
  \code{type_intrinsic} & \code{vfloat64m1_t}, \code{vint32mf2_t}, \dots; short arrays of
    vectors become RVV tuple types (\code{vfloat32m1x4_t}), because sizeless types
    cannot be struct members\\
  \code{constant_intrinsic} & uniform constants: \code{__riscv_vfmv_v_f_f64m1(c, 4uLL)};
    others (lane offsets $0,1,2,3$ of a vectorized counter):
    \code{pystencils::runtime::rvv::set_i64m1(0, 1, 2, 3)}\\
  \code{op_intrinsic} & \code{__riscv_vf{add,sub,mul,div}_vv_<sfx>} (integer:
    \code{v{add,sub,mul,div}}), broadcasts, casts, horizontal reductions\\
  \code{math_func_intrinsic} & \code{vfsqrt_v}, \code{vfabs_v}, \code{vfmin_vv},
    \code{vfmax_vv}; integer abs as \code{vmax(x, vneg(x))}; transcendentals and
    floor/ceil raise \code{MaterializationError}\\
  \code{vector_load}/\code{_store} & contiguous \code{__riscv_vle64_v_f64m1(&p[i], vl)}
    and \code{vse64}; strided \code{vlse64}/\code{vsse64} with a byte stride (x86 needs
    gathers and AVX-512 scatters here)\\
  \code{rng_engine_intrinsic}, \code{subscript_intrinsic} & vectorized Philox via the
    runtime header; tuple elements via \code{__riscv_vget_v_...}\\
  \bottomrule
\end{tabularx}
\end{center}
Casts use one intrinsic each, because the legal widths (128 and 256 bits at the same
lane count) differ by at most a factor of two: \code{vfcvt}, widening \code{vfwcvt} or
\code{vsext}, narrowing \code{vfncvt} or \code{vncvt}. Float-to-integer conversions
always use the \code{rtz} (round-towards-zero) variant, because C truncates.

For the waLBerla LBM sweep only a small subset matters: the loads and stores
\code{vle64} and \code{vse64}, the arithmetic \code{vfadd}, \code{vfsub}, \code{vfmul} and
\code{vfdiv}, and \code{vfmv_v_f} broadcasts (\cref{r02:lst:kernel}).

\section{The runtime header}\label{r02:sec:ps-runtime}

Generated code includes \file{pystencils_runtime/rvv.hpp}. The CMake integration of
pystencils-sfg puts \ps{}' include directory on the target's include path; on the board
that is \file{submodules/pystencils/src/pystencils/include} of the fork, found through
\code{PYTHONPATH}. It provides what the x86
backend takes from its own headers:
\begin{itemize}
  \item the VLEN guard (\cref{r02:sec:design-guard});
  \item \code{set_<sfx>(v0, v1, \dots)}: build a vector from per-lane values through a
        stack array and one load (RVV has no ``set'' intrinsic); generated by one macro
        for every legal type;
  \item \code{horizontal_{add,mul,min,max}_<sfx>(dst, v)}: store the lanes and fold them
        with a scalar loop, so the order of operations (and hence the result) matches
        scalar code exactly. Reductions run once per kernel call;
  \item \code{random::philox_float32/64}: wrappers around the shared vectorized Philox
        generator in \file{bits/philox_rand.h}, returning RVV tuple types.
\end{itemize}
The sweeps of \cref{r02:ch:results} use none of the helpers; for them the header is
only the guard and \code{<riscv_vector.h>}.

\section{JIT compiler flags}\label{r02:sec:ps-jit}

waLBerla generates code without compiling it in Python (\code{cfg.jit = no_jit}), so
this part matters for \ps{}' own tests and for Python experiments, not for the
waLBerla build. It mirrors waLBerla's CMake rule (\cref{r02:sec:wlb-cmake}):
\rneedspace{12\baselineskip}
\begin{pycode}[nolol=false,caption={[JIT compiler flags for RISC-V]Compiler flags in \texttt{jit/cpu/compiler\_info.py} (abridged).},label=r02:lst:jitflags]
_RISCV_RVV_MARCH = "-march=rv64gcv_zvl256b"

def _current_cpu_march() -> list[str]:
    """GCC (at least up to 14) rejects -march=native on RISC-V."""
    if platform.machine() == "riscv64":
        if Target.RISCV_RVV in Target.available_vector_cpu_targets():
            return [_RISCV_RVV_MARCH]
        return ["-march=rv64gc"]
    return ["-march=native"]

        match self.target:
            case Target.CurrentCPU:
                flags += _current_cpu_march()
            ...
            case Target.RISCV_RVV:
                flags += [_RISCV_RVV_MARCH]
\end{pycode}

\section{Tests}\label{r02:sec:ps-tests}

\begin{description}[style=nextline]
  \item[\file{tests/nbackend/test_rvv_target.py}] target flags, default lanes, ISA string
    parsing (\code{v} versus \code{zve*} and \code{xtheadvector}), JIT flags.
  \item[\file{tests/nbackend/test_rvv_codegen.py}] runs on any machine: the type and LMUL
    table, rejection of illegal types, intrinsic names for every construct, and a
    syntax-only compile of eight representative kernels with a riscv64 cross compiler
    (\code{riscv64-linux-gnu-g++ -march=rv64gcv_zvl256b}) or the native compiler on the
    board.
  \item[\file{tests/nbackend/test_rvv_execution.py}] runs only on RVV hardware: results
    against NumPy and \code{GenericCPU}, bit-exact Philox, truncating casts
    (\code{(int)-1.5 == -1}).
  \item[Shared modules] \file{test_vectorization.py}, \file{test_cpu_axes.py},
    \file{test_type_cast.py}, \file{test_functions.py}, \file{test_domain_kernels.py}
    and \file{test_flowgraph_codegen.py} got RVV cases next to the AVX and SVE ones.
\end{description}
Results are in \cref{r02:sec:results-tests}.
